% Options for packages loaded elsewhere
\PassOptionsToPackage{unicode}{hyperref}
\PassOptionsToPackage{hyphens}{url}
\PassOptionsToPackage{dvipsnames,svgnames,x11names}{xcolor}
\documentclass[
  10pt,
]{extarticle}
\usepackage{xcolor}
\usepackage[top=1.5in, bottom=1.5in, left=1.5in, right=1.5in]{geometry}
\usepackage{amsmath,amssymb}
\setcounter{secnumdepth}{5}
\usepackage{iftex}
\ifPDFTeX
  \usepackage[T1]{fontenc}
  \usepackage[utf8]{inputenc}
  \usepackage{textcomp} % provide euro and other symbols
\else % if luatex or xetex
  \usepackage{unicode-math} % this also loads fontspec
  \defaultfontfeatures{Scale=MatchLowercase}
  \defaultfontfeatures[\rmfamily]{Ligatures=TeX,Scale=1}
\fi
\usepackage{lmodern}
\ifPDFTeX\else
  % xetex/luatex font selection
\fi
% Use upquote if available, for straight quotes in verbatim environments
\IfFileExists{upquote.sty}{\usepackage{upquote}}{}
\IfFileExists{microtype.sty}{% use microtype if available
  \usepackage[]{microtype}
  \UseMicrotypeSet[protrusion]{basicmath} % disable protrusion for tt fonts
}{}
\makeatletter
\@ifundefined{KOMAClassName}{% if non-KOMA class
  \IfFileExists{parskip.sty}{%
    \usepackage{parskip}
  }{% else
    \setlength{\parindent}{0pt}
    \setlength{\parskip}{6pt plus 2pt minus 1pt}}
}{% if KOMA class
  \KOMAoptions{parskip=half}}
\makeatother
\usepackage{longtable,booktabs,array}
\newcounter{none} % for unnumbered tables
\usepackage{calc} % for calculating minipage widths
% Correct order of tables after \paragraph or \subparagraph
\usepackage{etoolbox}
\makeatletter
\patchcmd\longtable{\par}{\if@noskipsec\mbox{}\fi\par}{}{}
\makeatother
% Allow footnotes in longtable head/foot
\IfFileExists{footnotehyper.sty}{\usepackage{footnotehyper}}{\usepackage{footnote}}
\makesavenoteenv{longtable}
\usepackage{graphicx}
\makeatletter
\newsavebox\pandoc@box
\newcommand*\pandocbounded[1]{% scales image to fit in text height/width
  \sbox\pandoc@box{#1}%
  \Gscale@div\@tempa{\textheight}{\dimexpr\ht\pandoc@box+\dp\pandoc@box\relax}%
  \Gscale@div\@tempb{\linewidth}{\wd\pandoc@box}%
  \ifdim\@tempb\p@<\@tempa\p@\let\@tempa\@tempb\fi% select the smaller of both
  \ifdim\@tempa\p@<\p@\scalebox{\@tempa}{\usebox\pandoc@box}%
  \else\usebox{\pandoc@box}%
  \fi%
}
% Set default figure placement to htbp
\def\fps@figure{htbp}
\makeatother
\setlength{\emergencystretch}{3em} % prevent overfull lines
\providecommand{\tightlist}{%
  \setlength{\itemsep}{0pt}\setlength{\parskip}{0pt}}
\usepackage{xcolor}
\usepackage{hyperref}
\hypersetup{colorlinks=true, linkcolor=blue, urlcolor=blue, citecolor=blue}
\usepackage{float}
\usepackage{fvextra}
\usepackage{fancyhdr}
\usepackage{lastpage}
\usepackage{soul}
\usepackage{etoolbox}
\usepackage{microtype}
\usepackage{tcolorbox}
\usepackage{enumitem}
\usepackage{fancyvrb}
\usepackage{xspace}
\allowdisplaybreaks
\setcounter{tocdepth}{2}
\setcounter{secnumdepth}{0}
\floatplacement{figure}{H}
\makeatletter
\newcommand{\@subtitle}{}
\newcommand{\subtitle}[1]{\gdef\@subtitle{#1}}
\newcommand{\DocSubtitle}{\@subtitle}
\renewcommand{\maketitle}{
  \thispagestyle{plain}
  \null
  \vfill
  \begin{center}
    {\Large \textsc{\@title} \par}\vskip 0.5em
    {\LARGE \bfseries \@subtitle \par}\vskip 0.75em
    {\large \@author \par}\vskip 0.5em
    {\normalsize \@date \par}
  \end{center}
  \vfill
  \tableofcontents
  \vspace{1em}
  \clearpage
}
\AfterEndEnvironment{Highlighting}{\setcounter{CodeLine}{\numexpr\value{CodeLine}+\FV@CodeLineNo\relax}}
\makeatother
\fancypagestyle{plain}{
  \fancyhf{}
  \renewcommand{\headrulewidth}{0pt}
  \renewcommand{\footrulewidth}{0pt}
}
\pagestyle{fancy}
\fancyhf{}
\fancyfoot[C]{\small\textsc{Page}~\thepage~\textsc{of}~\pageref{LastPage}}
\fancyhead[L]{\textsc{Rento Saijo}}
\fancyhead[R]{\textsc{\DocSubtitle}}
\fancyhead[C]{\textsc{\rightmark}}
\renewcommand{\headrulewidth}{0.5pt}
\renewcommand{\footrulewidth}{0.5pt}
\newcounter{CodeLine}
\newcounter{cell}
\AtBeginEnvironment{Highlighting}{\stepcounter{cell}}
\DefineVerbatimEnvironment{Highlighting}{Verbatim}{
  frame=single,
  rulecolor=\color{black},
  framesep=2mm,
  label=\footnotesize Cell \thecell,
  labelposition=topline,
  commandchars=\\\{\},
  breaklines=true,
  breakanywhere=false,
  breaksymbol={},
  numbers=left,
  numbersep=3pt,
  firstnumber=\numexpr\value{CodeLine}+1\relax,
  fontsize=\small
}
\DefineVerbatimEnvironment{verbatim}{Verbatim}{
  breaklines=true,
  breakanywhere=true,
  numbers=left,
  numbersep=3pt,
  firstnumber=1,
  fontsize=\footnotesize
}
\newtcolorbox{problem}[1][]{
  title={\textsc{#1}}
}
\renewcommand{\contentsname}{Table of Contents}
\newcommand{\E}{\mathrm{E}}
\newcommand{\Var}{\mathrm{Var}}
\newcommand{\Cov}{\mathrm{Cov}}
\newcommand{\R}{\textsf{R}\xspace}
\newcommand{\SD}{\mathrm{SD}}
\newcommand{\SE}{\mathrm{SE}}
\usepackage{bookmark}
\IfFileExists{xurl.sty}{\usepackage{xurl}}{} % add URL line breaks if available
\urlstyle{same}
\hypersetup{
  pdftitle={STA 209: Introduction to Time Series Analysis},
  colorlinks=true,
  linkcolor={blue},
  filecolor={Maroon},
  citecolor={blue},
  urlcolor={blue},
  pdfcreator={LaTeX via pandoc}}

\title{STA 209: Introduction to Time Series Analysis}
\usepackage{etoolbox}
\makeatletter
\providecommand{\subtitle}[1]{% add subtitle to \maketitle
  \apptocmd{\@title}{\par {\large #1 \par}}{}{}
}
\makeatother
\subtitle{Extra Credit}
\author{Rento Saijo}
\date{May 12, 2026}

\begin{document}
\maketitle

\section{Disclosure}\label{disclosure}

GPT-5.4 (High) on Codex was used to create the \texttt{YAML} portion and some \texttt{LaTeX} code to format the text and equations. Page-formatting code was also provided by Derin Gezgin.

\newpage

\section{Problem 1}\label{problem-1}

\begin{problem}[Problem 1]
Consider a time series defined by
\[
X_t = \delta + X_{t-1} + W_t,
\]
for \(t=1,2,\ldots\) and \(X_0=0\), where \(W_t\) is a white noise with mean \(0\) and variance \(1\).
\end{problem}

\newpage

\subsection{Problem 1 (a)}\label{problem-1-a}

\begin{problem}[Problem 1 (a)]
Identify the model.
\end{problem}

The model
\[
X_t = \delta + X_{t-1} + W_t
\]
is a random walk with drift \(\delta\). That is, each new value equals the previous value, plus a constant drift term \(\delta\), plus a white-noise term \(W_t\).

\newpage

\subsection{Problem 1 (b)}\label{problem-1-b}

\begin{problem}[Problem 1 (b)]
Show that \(X_t\) can be written as \(X_t=t\delta+\sum_{j=1}^{t}W_j\).
\end{problem}

Starting from
\[
X_t = \delta + X_{t-1} + W_t,
\]
substitute recursively:
\[
\begin{aligned}
X_t
&= \delta + X_{t-1} + W_t \\
&= \delta + (\delta + X_{t-2} + W_{t-1}) + W_t \\
&= 2\delta + X_{t-2} + W_{t-1} + W_t \\
&= 2\delta + (\delta + X_{t-3} + W_{t-2}) + W_{t-1} + W_t \\
&= 3\delta + X_{t-3} + W_{t-2} + W_{t-1} + W_t \\
&\ \ \vdots \\
&= t\delta + X_0 + \sum_{j=1}^{t} W_j.
\end{aligned}
\]
Since the problem gives \(X_0=0\), this becomes
\[
X_t = t\delta + \sum_{j=1}^{t} W_j.
\]

\newpage

\subsection{Problem 1 (c)}\label{problem-1-c}

\begin{problem}[Problem 1 (c)]
Find the mean and ACVF of \(X_t\).
\end{problem}

From part (b),
\[
X_t = t\delta + \sum_{j=1}^{t} W_j.
\]

First, compute the mean:
\[
\begin{aligned}
\mathrm{E}(X_t)
&= \mathrm{E}\left(t\delta + \sum_{j=1}^{t} W_j\right) \\
&= \mathrm{E}(t\delta) + \sum_{j=1}^{t} \mathrm{E}(W_j) \\
&= t\delta + \sum_{j=1}^{t} 0 \\
&= t\delta.
\end{aligned}
\]
Hence,
\[
\mu_X(t) = t\delta.
\]

Now, compute the auto-covariance function. For integers \(h\) such that \(t+h \ge 0\),
\[
\gamma_X(t,h) = \mathrm{Cov}(X_{t+h},X_t) = \mathrm{E}[(X_{t+h}-\mu_{t+h})(X_t-\mu_t)].
\]
Since
\[
\mu_{t+h} = (t+h)\delta
\quad \text{and} \quad
\mu_t = t\delta,
\]
and from part (b),
\[
X_{t+h} = (t+h)\delta + \sum_{j=1}^{t+h} W_j
\quad \text{and} \quad
X_t = t\delta + \sum_{k=1}^{t} W_k,
\]
it follows that
\[
\begin{aligned}
X_{t+h}-\mu_{t+h}
&= \left((t+h)\delta + \sum_{j=1}^{t+h} W_j\right) - (t+h)\delta \\
&= \sum_{j=1}^{t+h} W_j,
\end{aligned}
\]
and
\[
\begin{aligned}
X_t-\mu_t
&= \left(t\delta + \sum_{k=1}^{t} W_k\right) - t\delta \\
&= \sum_{k=1}^{t} W_k.
\end{aligned}
\]
Therefore,
\[
\begin{aligned}
\gamma_X(t,h)
&= \mathrm{E}\left[\left(\sum_{j=1}^{t+h} W_j\right)\left(\sum_{k=1}^{t} W_k\right)\right] \\
&= \sum_{j=1}^{t+h} \sum_{k=1}^{t} \mathrm{E}(W_jW_k).
\end{aligned}
\]
Because \(W_t\) is white noise with variance \(1\),
\[
\mathrm{E}(W_jW_k) =
\begin{cases}
1, & j=k,\\
0, & j \neq k.
\end{cases}
\]

Now, let us evaluate the expression case-by-case. If \(h \ge 0\), then
\[
\begin{aligned}
\gamma_X(t,h)
&= \sum_{j=1}^{t+h} \sum_{k=1}^{t} \mathrm{E}(W_jW_k).
\end{aligned}
\]
Here, the only matching indices are
\[
j=k=1,2,\ldots,t,
\]
so there are \(t\) nonzero terms, each equal to \(1\). Therefore,
\[
\gamma_X(t,h)=t.
\]

If \(-t \le h < 0\), then \(t+h < t\), so the only matching indices are
\[
j=k=1,2,\ldots,t+h.
\]
Thus, there are \(t+h\) nonzero terms, each equal to \(1\). Therefore,
\[
\gamma_X(t,h)=t+h.
\]

Hence, the ACVF is
\[
\gamma_X(t,h) =
\begin{cases}
t, & h \ge 0,\\
t+h, & -t \le h < 0.
\end{cases}
\]

\newpage

\subsection{Problem 1 (d)}\label{problem-1-d}

\begin{problem}[Problem 1 (d)]
Is this time series stationary? Justify.
\end{problem}

No, \(X_t\) is not stationary. For (weak) stationarity, two conditions are required:

\begin{enumerate}
\def\labelenumi{\arabic{enumi}.}
\tightlist
\item
  The mean function must not depend on time \(t\).
\item
  The ACVF must depend only on lag \(h\), not on time \(t\).
\end{enumerate}

But here,
\[
\mu_X(t) = t\delta,
\]
which depends on \(t\) unless \(\delta=0\). So, in general, the mean is not constant over time. Therefore, the process already fails the first condition for weak stationarity. Hence, \(X_t\) is not stationary.

\end{document}
